\chapter*{Resumen}
\addcontentsline{toc}{chapter}{Resumen}
\markboth{RESUMEN}{RESUMEN}

\begin{singlespace}
Este trabajo estudia la cotización de una renta vitalicia inmediata en un modelo
particular que incorpora piso pensional, indexación y auxilio funerario. La mesada
inicial $M$ se obtiene al resolver $f(M)=\VP(M)-\KA=0$, donde $\VP$ es el valor
presente actuarial de las obligaciones incluidas y $\KA$ el capital objetivo,
expresados sobre una misma base de valoración. La formulación estudia el cotizador
implementado y no constituye una representación exhaustiva de la regulación ni de
la tarificación comercial de las aseguradoras colombianas.

Bajo las hipótesis del modelo, $\VP$ es \emph{lineal a trozos}, con a lo sumo
$\datoNQuiebresMax$ umbrales candidatos agrupados en tres racimos geométricos
alrededor de $1$, $5$ y $10$ salarios mínimos para la parametrización considerada.
La extensión a $M\ge0$ es constante en $[0,\SMLV]$; en el dominio admisible
$M\ge\SMLV$, existe una solución única si y solo si $\KA\ge\VP(\SMLV)$.
El auxilio funerario puede producir pérdida de convexidad. Se obtiene una
expresión cerrada de las derivadas laterales y del $B$\nobreakdash-subdiferencial,
con interpretación como factores marginales de renta actuarial. Newton con una
pendiente exacta alcanza la raíz en un paso cuando el tramo seleccionado contiene
también la solución en su clausura; no se identifica el número de pasos con el de
tramos atravesados. Una salvaguarda de encaje con bisección periódica proporciona
convergencia global bajo las condiciones establecidas.

Se comparan \datoNMetodos{} métodos en \datoNEscenarios{} escenarios, que incluyen
\datoNSinteticos{} sintéticos, \datoNReales{} cotizaciones de renta vitalicia y
\datoNRealesRP{} registros adicionales con información de retiro programado.
La batería utiliza el objetivo RV; las trayectorias anuales de RP se estudian
separadamente. La comparación distingue
iteraciones, llamadas efectivas al motor y tiempo de ejecución; contabiliza la
preparación y utiliza una evaluación conjunta de valor y derivadas cuando estas
son analíticas. Los resultados se interpretan sobre los escenarios examinados y
no como una clasificación universal de los métodos.

El retiro programado se formula mediante el estado acoplado de capital y mesada
anterior. Una cota condicional de propagación permite relacionar la tolerancia
anual con un límite de error del saldo final. Se examinan además las fuentes y
limitaciones de las variables actuariales, financieras y de calibración, y se
documenta la implementación en Python. La precisión numérica y la concordancia
con referencias de calibración no sustituyen la validación actuarial del modelo.
\end{singlespace}

\vspace{.5cm}
\noindent\textbf{Palabras clave:} ecuaciones no lineales; métodos iterativos;
análisis no suave; Newton semisuave; funciones lineales a trozos; matemática
actuarial; rentas vitalicias; retiro programado.

\vspace{.4cm}
%\noindent\textbf{Clasificación MSC 2020:} 65H05, 49J52, 91G05, 90C31.
